% Source: https://artofproblemsolving.com/community/c3h2142190_hard_mock_amc_10_problem
% Stable identity: aops-topic-2142190-q01
Since $AB^2=212$, the radius is $\dfrac{AB}{2}$ with $\left(\dfrac{AB}{2}\right)^2=\dfrac{212}{4}=53$, so the circle is $x^2+y^2=53$.

\textbf{Point $E$.} From $13x+5y=9$, $y=\dfrac{9-13x}{5}$. Substituting,
\[25x^2+(9-13x)^2=1325\iff194x^2-234x-1244=0\iff97x^2-117x-622=0\iff(x+2)(97x-311)=0 .\]
The root $x=-2$ gives $y=7$, the point $(-2,7)$ in quadrant II; the other root gives $F=\left(\dfrac{311}{97},-\dfrac{634}{97}\right)$ in quadrant IV. So $E=(-2,7)$.

\textbf{Points $G$ and $H$.} Translating $y=x$ five units to the left gives $y=x+5$. Then
\[x^2+(x+5)^2=53\iff x^2+5x-14=0\iff(x-2)(x+7)=0,\]
so $G$ and $H$ are $(2,7)$ and $(-7,-2)$.

\begin{asy}
usepackage("amsmath");
size(6.5cm);
defaultpen(fontsize(10pt));
real R=sqrt(53);
pair Ep=(-2,7), F=(311/97,-634/97), G=(2,7), H=(-7,-2);
filldraw(Ep--G--H--cycle,paleyellow,linewidth(0.9));
draw((-9,0)--(9,0),gray,Arrow(4)); draw((0,-9)--(0,9),gray,Arrow(4));
label("$x$",(9,0),S,gray); label("$y$",(0,9),W,gray);
draw(circle((0,0),R),linewidth(0.9));
draw((-3,9.6)--(4,-8.6),blue+linewidth(0.8));
draw((-9,-4)--(3.5,8.5),red+linewidth(0.8));
draw(Ep--(-2,-2),dashed);
dot(Ep); dot(F,blue); dot(G); dot(H); dot((0,0));
label("$E(-2,7)$",Ep,NW); label("$G(2,7)$",G,SE); label("$H(-7,-2)$",H,W);
label("$F$",F,E,blue); label("$O$",(0,0),SE);
label("$13x+5y=9$",(-3,9.6),W,blue+fontsize(8pt)); label("$y=x+5$",(-9,-4),S,red+fontsize(8pt));
label("\mbox{}",(min(currentpicture,user=true)+max(currentpicture,user=true))/2,black); // keeps the text after the figure black when LaTeX typesets the labels inline
\end{asy}

\textbf{Area.} The points $E=(-2,7)$ and $G=(2,7)$ lie on the line $y=7$, so $EG=4$, and $H=(-7,-2)$ is $7-(-2)=9$ below that line. The area is
\[\dfrac{1}{2}\cdot4\cdot9=\boxed{18}.\]
